\documentclass[11pt]{article}

\usepackage[a4paper,margin=1in]{geometry}
\usepackage{booktabs}
\usepackage{amsmath}
\usepackage{xcolor}
\usepackage{listings}
\usepackage{enumitem}
\usepackage[hidelinks]{hyperref}

\definecolor{codebg}{rgb}{0.96,0.96,0.96}
\lstset{
  basicstyle=\ttfamily\small,
  backgroundcolor=\color{codebg},
  breaklines=true,
  frame=single,
  columns=fullflexible,
  keepspaces=true,
  showstringspaces=false,
}

\newcommand{\code}[1]{\texttt{#1}}

\title{\textbf{Fine-Tuning Qwen2.5-Coder-14B for\\ Natural-Language $\rightarrow$ GeoGebra Geometry Construction}\\[4pt]
\large Training Report (Kaggle, QLoRA)}
\author{Aditya Dhull}
\date{June 25, 2026}

\begin{document}
\maketitle

\begin{abstract}
This report documents the fine-tuning of \textbf{Qwen2.5-Coder-14B-Instruct} to translate
natural-language Olympiad-geometry problems into ordered \textbf{GeoGebra} construction
programs encoded as JSON, which render as interactive diagrams. The model was trained with
\textbf{QLoRA} on a free \textbf{Kaggle} GPU (NVIDIA T4, 16~GB) over a curated and
paraphrase-augmented dataset of geometry problems. Training completed in roughly two hours.
On 31 held-out problems the model produced \textbf{100\% syntactically valid JSON} and
\textbf{71\% fully renderable, non-degenerate figures} with no repair pass --- exceeding the
earlier 3B baseline (65\% \emph{with} a repair loop). We describe the dataset, the model and
hyperparameters, the training environment and procedure, the results and failure analysis,
the deployment constraints, and the steps to reproduce.
\end{abstract}

\section{Problem and Approach}
The system maps a natural-language geometry problem to an \emph{ordered} list of GeoGebra
construction steps (JSON), each step being a named object with a native GeoGebra command that
references previously defined names. Rendering the steps in a GeoGebra applet yields an
interactive, draggable diagram. The schema for one step is:

\begin{lstlisting}
{ "name": <unique id>, "type": "final" | "construction",
  "class": point|line|segment|ray|circle|conic|locus|number|angle,
  "label": <optional caption>, "command": <native GeoGebra syntax> }
\end{lstlisting}

\code{final} objects are shown; \code{construction} objects are hidden scaffolding. A
\emph{visibility rule} enforces that every entity named in the problem statement is marked
\code{final} (and thus visible).

A prior fine-tune of \textbf{Qwen2.5-Coder-3B} (on a 6~GB GPU) learned the output \emph{format}
but frequently emitted incorrect GeoGebra commands (e.g.\ \code{Incircle(I, line)} instead of
\code{Incircle(A,B,C)}). Because such command errors stem from the small base model's limited
pretrained knowledge rather than from a lack of data, we scaled the base model to 14B
parameters; this report covers that 14B run.

\section{Dataset}
\paragraph{Composition.} The corpus contains \textbf{232 distinct problems}: 82 hand-curated
``gold'' problems and 150 problems generated and then verified inside a real (headless)
GeoGebra engine for both renderability and the geometric property they assert.

\paragraph{Split.} A deterministic, source-stratified split (seed 42, $\sim$13\% eval) yields
\textbf{201 training} and \textbf{31 evaluation} problems (the eval set comprises 11 gold $+$
20 generated and is never trained on).

\paragraph{Augmentation.} The \emph{training} split is paraphrase-augmented: each statement is
restated in meaning-preserving ways (synonyms, openers, British/American spelling) while the
construction target is unchanged. This expands the training set $201 \rightarrow 658$ and
improves robustness to phrasing. The eval split is left original to keep the measurement
honest.

\paragraph{Format.} Each example is a chat triple --- a fixed \emph{system} prompt describing
the schema and the visibility rule, a \emph{user} turn (the problem statement), and an
\emph{assistant} turn (the JSON construction). The longest example is $\approx$1{,}295 tokens.

\begin{table}[h]
\centering
\caption{Dataset composition and split.}
\begin{tabular}{lrr}
\toprule
 & \textbf{Train} & \textbf{Eval} \\
\midrule
Gold (hand-curated)      & 71  & 11 \\
Generated (validated)    & 130 & 20 \\
\midrule
Distinct problems        & 201 & 31 \\
After paraphrase aug.    & \textbf{658} & 31 (unchanged) \\
\bottomrule
\end{tabular}
\end{table}

\section{Model and Method}
\paragraph{Base model.} \code{Qwen2.5-Coder-14B-Instruct}, loaded in 4-bit
(\code{unsloth/Qwen2.5-Coder-14B-Instruct-bnb-4bit}), chosen for its strong code/structured-output
ability and permissive Apache-2.0 license.

\paragraph{Method.} \textbf{QLoRA}: low-rank adapters trained on top of the frozen 4-bit base,
via the Unsloth library (which provides the memory optimizations that let a 14B fit a single
16~GB GPU). Only the response (assistant JSON) contributes to the loss; the prompt is masked.

\paragraph{Trainable parameters.} \textbf{137{,}625{,}600} of \textbf{14{,}907{,}659{,}264}
($\approx$0.92\%).

\begin{table}[h]
\centering
\caption{Training hyperparameters.}
\begin{tabular}{ll}
\toprule
\textbf{Setting} & \textbf{Value} \\
\midrule
LoRA rank $r$ / $\alpha$ / dropout & 32 / 32 / 0 \\
LoRA target modules & q,\,k,\,v,\,o,\,gate,\,up,\,down \\
Max sequence length & 1536 \\
Per-device batch size & 1 \\
Gradient accumulation & 8 (effective batch 8) \\
Epochs & 2 \\
Total optimizer steps & 166 \\
Learning rate / schedule & $2\times10^{-4}$ / cosine \\
Warmup ratio & 0.05 \\
Optimizer & adamw\_8bit \\
Weight decay & 0.01 \\
Precision & fp16 (T4 has no bf16) \\
Gradient checkpointing & enabled (Unsloth) \\
Loss masking & response-only (ChatML markers) \\
\bottomrule
\end{tabular}
\end{table}

\paragraph{Epoch choice.} Because the augmentation makes each construction appear $\sim$3 times
per epoch, 2 epochs ($\sim$6 exposures per construction) is the sweet spot: it halves runtime
versus 4 epochs \emph{and} reduces overfitting on the 232 distinct constructions.

\section{Training Environment and Procedure}
Training ran on a Kaggle notebook with a single GPU (Unsloth is single-GPU; the second T4 in a
``T4 $\times$2'' instance cannot accelerate a 14B that already needs Unsloth's memory savings
to fit one 16~GB card).

\begin{table}[h]
\centering
\caption{Software / hardware environment.}
\begin{tabular}{ll}
\toprule
GPU & NVIDIA Tesla T4, 14.56~GB usable, compute 7.5 \\
Peak GPU memory & $\approx$12.7~GB \\
PyTorch / CUDA & 2.10.0+cu128 / 12.8 \\
Triton & 3.6.0 \\
Transformers & 5.5.0 \\
Unsloth & 2026.6.9 \\
Attention backend & xFormers 0.0.35 (FA2 unavailable) \\
\bottomrule
\end{tabular}
\end{table}

\paragraph{Runtime and loss.} Training took \textbf{6{,}917~s} ($\approx$1~h~55~m) for 166
steps ($\approx$32--42~s/step). The mean training loss was \textbf{0.0346}; per-step loss fell
smoothly from $\sim$0.16 to $\sim$0.006.

\begin{table}[h]
\centering
\caption{Representative training-loss trajectory.}
\begin{tabular}{rr}
\toprule
\textbf{Epoch} & \textbf{Loss} \\
\midrule
0.12 & 0.1606 \\
0.49 & 0.0578 \\
0.85 & 0.0254 \\
1.09 & 0.0143 \\
1.57 & 0.0078 \\
1.94 & 0.0070 \\
\bottomrule
\end{tabular}
\end{table}

\section{Results}
The trained model was evaluated on the 31 held-out problems by generating a construction for
each (greedy decoding, up to 1536 new tokens), then rendering every construction in a headless
GeoGebra engine across 7 random parameterizations. A construction counts as ``valid'' only if
\emph{all} trials render with no broken steps and no degeneracy (points at infinity, collapse,
mass coincidence).

\begin{table}[h]
\centering
\caption{Functional evaluation on 31 held-out problems (14B vs.\ 3B baseline; identical split).}
\begin{tabular}{lcc}
\toprule
\textbf{Model} & \textbf{Parseable JSON} & \textbf{Renders \& valid} \\
\midrule
Qwen2.5-Coder-3B  & 100\% & 61\% raw / \textbf{65\% with repair loop} \\
Qwen2.5-Coder-14B & 100\% & \textbf{71\% raw (no repair)} \\
\bottomrule
\end{tabular}
\end{table}

The 14B at \textbf{71\% raw} surpasses the 3B's \textbf{65\% with} a repair pass; the repair
loop was not applied to the 14B. On simple single-construction prompts (incircle and contact
triangle, the three altitudes, two tangents from an external point) the 14B produced correct,
complete figures.

\paragraph{Failure analysis (9 of 31).} The remaining failures are command-level issues:
\begin{itemize}[noitemsep]
  \item \textbf{Hallucinated commands} (plausible names that are not GeoGebra): \code{Incenter(...)},
        \code{PerpendicularFoot(...)}, \code{Reflection(...)} (the correct name is \code{Reflect}).
  \item \textbf{Malformed commands}: \code{Intersect(Circle(O2,P1), 1)} (missing second object),
        \code{AngleBisector(A, B)} (needs three points), \code{Intersect(incircle, seg)} (missing index).
  \item \textbf{Undefined references}: e.g.\ \code{Line(A, Aprime)} with \code{Aprime} never defined.
\end{itemize}
Most of these are mechanical and are exactly what the generate$\rightarrow$validate$\rightarrow$repair
loop corrects (feed the failing step and error back); applying it is expected to lift the rate
toward $\sim$80\%.

\section{Deployment}
Producing a single quantized GGUF for local serving requires merging the LoRA into a 16-bit
copy of the 14B and quantizing it, which peaks at $\approx$58~GB of disk --- exceeding Kaggle's
$\approx$57~GB, so the GGUF export fails with ``no space left on device.'' The training itself
is unaffected; the LoRA adapter ($\approx$290~MB) saves successfully. The model is therefore
served by \textbf{loading the 4-bit base $+$ adapter} in a Kaggle/Colab inference notebook
(no GGUF needed). Generated JSON is pasted into a single-file GeoGebra renderer to view the
interactive diagram.

\section{Reproducibility}
The pipeline is script-driven:
\begin{itemize}[noitemsep]
  \item \code{prepare\_data.py} --- merges the 232 problems, applies the visibility rule, augments
        the train split, and writes \code{train.jsonl} (658) / \code{eval.jsonl} (31).
  \item \code{augment.py} --- meaning-preserving paraphrasing of statements.
  \item \code{qlora\_train\_14b.py} --- the QLoRA training described above.
  \item Inference cells --- load the adapter and generate JSON for new problems.
\end{itemize}

A minimal inference recipe (Kaggle: GPU + Internet on, \code{geom14b-adapter} dataset attached):
\begin{lstlisting}
from unsloth import FastLanguageModel
import glob, json
ADAPTER = glob.glob("/kaggle/input/**/adapter_config.json", recursive=True)[0].rsplit("/",1)[0]
m, tok = FastLanguageModel.from_pretrained(ADAPTER, max_seq_length=2048, load_in_4bit=True)
FastLanguageModel.for_inference(m)
# build chat prompt (system schema + user problem), greedy generate, parse JSON, lowercase class
\end{lstlisting}

\section{Limitations and Future Work}
\begin{itemize}[noitemsep]
  \item \textbf{Data breadth.} Only 232 \emph{distinct} constructions; augmentation adds phrasing
        variety, not new construction types. More distinct, validated problems would widen coverage.
  \item \textbf{Command hallucinations.} A short ``use exactly these GeoGebra commands'' hint in the
        prompt, and/or more data, would suppress invented names like \code{Reflection}.
  \item \textbf{Repair loop.} Not yet wired into the 14B inference path; doing so should raise the
        usable rate to $\sim$80\%.
  \item \textbf{Logical correctness.} The validator checks \emph{renderability}, not whether the figure
        faithfully matches the statement; a machine-checkable goal field would close this gap.
  \item \textbf{Serving.} A free, always-on API for a 14B is impractical; on-demand notebook
        inference (or a smaller distilled model with a GGUF) is the realistic path.
\end{itemize}

\section{Conclusion}
Scaling the base model to Qwen2.5-Coder-14B and fine-tuning with QLoRA on a free Kaggle T4
produced a model that reliably emits valid, complete GeoGebra constructions for most held-out
problems (100\% parseable, 71\% renderable raw), clearly improving on the 3B baseline and
fixing the dominant command-error mode at little additional cost. The remaining gap is
addressable with the existing repair loop, a small prompt refinement, and additional distinct
training data.

\end{document}
